\documentclass[11pt,a4paper]{article}
\usepackage[margin=27mm]{geometry}
\usepackage{amsmath,amssymb,amsthm,hyperref}
\newtheorem{theorem}{Theorem}
\newtheorem{lemma}[theorem]{Lemma}
\newtheorem{corollary}[theorem]{Corollary}
\newcommand{\F}{\mathbb F}
\newcommand{\Sz}{\operatorname{Sz}}
\newcommand{\gen}[1]{\langle #1\rangle}
\title{Scalar subspaces in rank-two Suzuki geometries}
\author{Human Compression project}\date{}
\begin{document}\maketitle
\begin{abstract}
A maximal subfield need not be maximal as an invariant additive subspace.
This distinction changes a family in the rank-two Suzuki classification.
We give the exact overgroups in the relevant mixed parabolic, construct
the resulting geometries, and determine their equivalence by field
automorphisms. For $\Sz(32768)$ this single family has forty-one
isomorphism classes in the source's convention, whereas its table gives
fourteen classes in total. The argument reuses the scalar-center
correspondence already established in this project's work on higher-rank
Suzuki configurations. It does not classify the remaining families.
\end{abstract}

\section{An elementary overgroup correspondence}
We first state the reusable group-theoretic fact. All groups in this note
are finite. An action is called primitive when it is transitive and has
no nontrivial invariant partition.

\begin{lemma}\label{module}
Let $V$ be an abelian normal subgroup of $Q$, let $H\le Q$, and suppose
$Q=VH$. Put $W=V\cap H$. The overgroups of $H$ in $Q$ correspond
bijectively to the $H$-invariant subgroups of $V$ containing $W$, by
\[
             L\longmapsto L\cap V,\qquad T\longmapsto TH.
\]
The coset action on $Q/H$ is primitive precisely when $V/W$ has no
proper nonzero $H$-invariant subgroup. It is doubly transitive precisely
when $H$ is transitive on the nonzero elements of $V/W$.
\end{lemma}
\begin{proof}
If $H\le L\le Q$, write an element of $L$ as $vh$. Since $h\in L$,
we have $v\in L\cap V$, and therefore $L=(L\cap V)H$. Conversely,
an $H$-invariant subgroup $T\le V$ makes $TH$ a subgroup; because
$W\le T$, its intersection with $V$ is $T$. These operations are
inverse. The usual correspondence between blocks containing the base
coset and overgroups of its stabilizer proves the assertion about
primitivity. Finally, $V/W$ acts regularly on the cosets, identified
with its additive set, and the base stabilizer acts by conjugation.
Its transitivity away from zero is exactly double transitivity.
\end{proof}

If $V$ is a vector space over a field $K$ and the image of $H$ on $V$
is the full group of nonzero scalar multiplications, the invariant
additive subgroups are exactly the $K$-subspaces. Thus an overgroup
$TH$ covers $H$ precisely when $\dim_K(T/W)=1$. On those cosets its
permutation group is the affine group of the line over $K$.
This is the scalar-center argument from the earlier project manuscript;
the statement above isolates the group-theoretic reason for it.

\section{Application to the Suzuki root group}
Let $n>d\ge3$ be odd integers, with $d\mid n$, and put
$F=\F_{2^n}$, $K=\F_{2^d}$, $q=2^n$, $s=2^d$ and
$\theta=2^{(n+1)/2}$. We use the standard faithful four-dimensional
Suzuki matrices, with
\[
 U(a,b)=\begin{pmatrix}
 1&0&0&0\\ a&1&0&0\\ a^{\theta+1}+b&a^\theta&1&0\\
 a^{\theta+2}+ab+b^\theta&b&a&1
 \end{pmatrix}.
\]
Multiplication gives the root-group law
\[
 U(a,b)U(c,e)=U(a+c,b+e+ac^\theta),\qquad
 Z_F=\{U(0,b):b\in F\}.
\]
The diagonal torus consists of
\[
 h(t)=\operatorname{diag}(t^{\theta+1},t,t^{-1},t^{-\theta-1}),
 \qquad
 h(t)^{-1}U(a,b)h(t)=U(t^\theta a,t^{\theta+2}b).
                                                        \tag{1}
\]
Let $U_K$ be given by $a,b\in K$, let $H_K=\{h(t):t\in K^*\}$,
and put $B=U_KH_K$. These are the standard root group, torus and
point stabilizer in $S=\Sz(s)\le G=\Sz(q)$. The ambient exponent
$\theta$ restricts to the Suzuki exponent on $K$: the two exponent
indices differ by a multiple of $d$.

We retain the classical Suzuki facts that this matrix model is faithful,
that $B$ is the stabilizer in the doubly transitive action of $S$ on
$s^2+1$ ovoid points, that $S$ is maximal in $G$ when $n/d$ is prime,
and that $\operatorname{Aut}(G)=G\rtimes\gen{\phi}$ with
$\phi:x\mapsto x^2$ of order $n$. These structural results are due to
Suzuki and Tits, and are also the structural inputs to the source
classification: see its Section 2.2 and the structural list in Section 4,
which cite Suzuki's \emph{On a class of doubly transitive groups}
(Annals of Mathematics 75 (1962), 105--145), Tits and L\"uneburg.
We do not reprove these classical facts here.

Both exponents in (1) are relatively prime to $q-1$. This is immediate
for $\theta$. For the other, an odd common divisor of
$2^n-1$ and $2^{(n-1)/2}+1$ would give simultaneously
$2^n=1$ and $2^{n-1}=1$ modulo that divisor, forcing it to be one.
The action of $H_K$ on $Z_F$ is consequently the full group of
$K$-scalar multiplications.

\begin{theorem}\label{covers}
Put $Q=Z_FB$. The overgroups of $B$ in $Q$ are exactly
\[
                   Q_T=Z_TB\quad (K\le T\le F),           \tag{2}
\]
where $T$ ranges over $K$-subspaces and $Z_T=\{U(0,t):t\in T\}$.
Their orders are $|Q_T|=|T|s(s-1)$. In particular, the groups covering
$B$ are exactly those with $\dim_K T=2$, and their order is
$s^3(s-1)$.
For every $c\in F\setminus K$,
\[
                   \gen{B,U(0,c)}=Z_{K+Kc}B.              \tag{3}
\]
\end{theorem}
\begin{proof}
Here $Z_F$ is abelian and normal in $Q$, and $Z_F\cap B=Z_K$.
The root group $U_K$ centralizes $Z_F$; the diagonal action is the
scalar action just proved. Lemma~\ref{module} gives (2), the order
formula, and the maximality criterion. Conjugating $U(0,c)$ by $H_K$
generates exactly the additive space $Kc$, and $B$ already contains
$Z_K$. This proves (3).
\end{proof}

The source's Lemma 6.9, printed page 112, asserts instead that (3) has
order $qs(s-1)$ and contains $B$ maximally because $K$ is a maximal
subfield. If $n/d=r$ is an odd prime, then $r\ge3$. There are
$K$-subspaces strictly between $K$ and $F$, regardless of the absence
of intermediate fields. Thus $B$ is never maximal in this $Q$.
This invalidates the corresponding mixed-parabolic row of Tables 8
and 11. The distinction is visible already at $s=8$, $q=512$:
\[
                   |B|=448<3584<|Q|=28672.                \tag{4}
\]

\section{The actual geometries and their equivalence}
Suppose from now on that $r=n/d$ is prime. For every two-dimensional
$K$-subspace $T$ containing $K$, form the rank-two coset geometry
\[
                     \Gamma_T=\Gamma(G;S,Q_T).            \tag{5}
\]
As in the source, elements are left cosets and incidence means
nonempty intersection. Isomorphism means an automorphism of the
specified group $G$ carrying one subgroup configuration to the other,
with type interchange allowed. We make no assertion here about
isomorphisms that do not preserve the specified group action.

\begin{theorem}\label{geometries}
Every geometry (5) is firm, residually connected, flag transitive,
residually weakly primitive, and locally doubly transitive. Its two
local degrees are $s^2+1$ and $s$.
Two members are conjugate under $G$ precisely when their spaces $T$
are equal. Their isomorphism classes in the source's convention are
exactly the orbits of $\phi$ on
\[
                        \mathbb P_K(F/K).                 \tag{6}
\]
There are $(s^{r-1}-1)/(s-1)$ conjugacy classes before this field fusion.
\end{theorem}
\begin{proof}
Coordinates in the standard subfield matrices give
$S\cap Q_T=B$. More explicitly, every element of $Q_T$ has a unique
expression $U(a,b)h(t)$ with $a\in K$, $b\in T$, $t\in K^*$;
such an element is in $S$ exactly when $b\in K$.
The subgroup $S$ is maximal in $G$. Since $T\ne K$, the group $Q_T$
is not contained in $S$, and therefore $\gen{S,Q_T}=G$.
The incidence graph is consequently connected. A rank-two coset
system is a geometry and is flag transitive: every element has an
incident partner, and each incident pair is a translate of the base
pair. Its two valencies are $[S:B]=s^2+1$ and $[Q_T:B]=s$,
so it is firm. In rank two, connectedness supplies all the residual
connectivity conditions of rank at least two.

The stabilizer $B$ is maximal in $S$ by the classical primitive ovoid
action, and is maximal in $Q_T$ by Theorem~\ref{covers}. Along with
maximality of $S$ in $G$, these are exactly the source's rank-two
residual weak primitivity conditions. The action of $S$ on its local
cosets is doubly transitive. By Lemma~\ref{module}, the other local
action is the affine action of $K\rtimes K^*$ on $K$, also doubly
transitive. This proves the first assertions.

The group $S$ is nonsoluble whereas $Q_T$ is soluble, so an
isomorphism cannot exchange their types. The joint normalizer of
$S$ and $B$ in $G$ is $B$: maximality and simplicity give
$N_G(S)=S$, and $B$ is self-normalizing in $S$ because it fixes a
unique ovoid point. Every element of $B$ preserves each $T$, by
the scalar action. Hence no two different spaces in (2) give
conjugate configurations.

Both $S$ and $B$ are preserved by $\phi$. Their joint normalizer
in $\operatorname{Aut}(G)$ is thus $B\rtimes\gen{\phi}$.
After the trivial $B$ action is removed, equivalence is exactly
field fusion of the spaces $T$. Such spaces are the inverse images
of lines in $F/K$, giving (6) and its stated cardinality.
\end{proof}

\section{Two exact field counts}
The following counts use the full automorphism group of $G$, not just
inner conjugacy. They require no Suzuki group census.

\begin{lemma}\label{fixed}
Let $\sigma$ be an odd-order field automorphism of $F$ preserving $K$.
Every $\sigma$-fixed element of $F/K$ has a representative in
$F^\sigma$. Consequently
\[
                       (F/K)^\sigma=F^\sigma/K^\sigma.
\]
\end{lemma}
\begin{proof}
If $\sigma$ has odd order $a$ and the class of $v$ is fixed, the sum
$v+\sigma(v)+\cdots+\sigma^{a-1}(v)$ is a fixed representative of
that same class: its image in $F/K$ is $av=v$ in characteristic two.
The kernel of the map from fixed representatives is $K^\sigma$.
\end{proof}

First take $K=\F_8$, $F=\F_{512}$. There are nine lines in $F/K$.
The map $\phi^3$ is $K$-linear of order three. An invariant line
would have an eigenvalue in $K^*$ of order dividing three, hence
eigenvalue one. But its fixed space is zero by Lemma~\ref{fixed},
since $F^{\phi^3}=K$. Therefore $\phi^3$ fixes no line. The orbits
of the order-nine map $\phi$ can only have lengths one, three or nine;
the first two would be fixed by $\phi^3$. All nine lines form one orbit.

Now take $K=\F_8$, $F=\F_{32768}$. There are
\[
                         (8^4-1)/7=585                  \tag{7}
\]
lines in $F/K$. The $K$-linear map $\phi^3$ has order five and no
invariant line: the same argument uses $\gcd(5,7)=1$ and
$F^{\phi^3}=K$. Hence any power $\phi^j$ with $5\nmid j$ fixes
no line, because its generated group contains $\phi^3$.

Put $\sigma=\phi^5$, of order three. Its fixed field is $\F_{32}$,
and its fixed field on $K$ is $\F_2$. Thus Lemma~\ref{fixed} gives
a fixed subspace of $F/K$ with sixteen elements. On $K$, $\sigma$
acts as the fourth power. If a $K$-line is invariant, choose a
nonzero vector $v$ on it and write $\sigma(v)=\lambda v$.
There is a unique $k\in K^*$ with $k^3=\lambda^{-1}$, because
$\gcd(3,7)=1$. Then $kv$ is fixed. Conversely a fixed nonzero vector
spans an invariant line. Each such line contains exactly one nonzero
fixed vector, since the fixed scalar field is $\F_2$. There are
therefore fifteen fixed lines for each of $\phi^5$ and $\phi^{10}$.
Burnside's orbit count now gives
\[
                            \frac{585+15+15}{15}=41.     \tag{8}
\]
There are three orbits of length five and thirty-eight of length
fifteen.

\begin{corollary}
For $G=\Sz(32768)$, the geometries (5) with $S=\Sz(8)$ give
forty-one isomorphism classes in the source's convention. All are
locally doubly transitive. Theorem 7.1 and its Table 11 omit this
family, and the total fourteen in Table 12 at $e=7$ is false.
\end{corollary}
\begin{proof}
The count is (8). The group $Q_T$ has order $3584$ and a nonabelian
normal Sylow two-subgroup. None of the other rows with the nonsoluble
parabolic $\Sz(8)$ has this subgroup type: their other parabolic is
Suzuki, dihedral, a cyclic group extended by a group of order four,
or the mixed group of order $32768\cdot8\cdot7$. Thus (5) is absent
from the asserted exhaustive list. Forty-one classes in one family
also contradict the printed total fourteen. The same order argument
already gives an omitted family at $q=512$.
\end{proof}

\paragraph{Attribution and scope.}
The target is D.~Leemans, \emph{The Rank 2 Geometries of the Simple
Suzuki Groups $\Sz(q)$}, Beitr\"age zur Algebra und Geometrie 39(1)
(1998), 97--120: Lemma 6.9, Theorem 7.1 and Tables 8, 11, 12.
The source's classification is considered under its specified-action
equivalence, defined on printed page 100. The matrix model, ovoid
action, subgroup maximality and ambient automorphism description are
retained classical inputs. The scalar-center argument was already
proved and separately reviewed in the project's earlier Suzuki
subfield manuscript, Section 4; the new contribution here is its
application to this rank-two row, the genuine replacement family,
its local double transitivity and its equivalence calculation.
No historical priority is asserted. The remaining rows and the
complete corrected global count are not established by this note.
This corrects a published classification; it is not an answer to a
named open conjecture.

The supplementary program independently constructs the fields of
orders $512$ and $32768$, verifies their defining polynomials by
Rabin's criterion, lists all two-dimensional $\F_8$-subspaces
containing $\F_8$, and checks every field orbit. It agrees with
(7)--(8). These finite controls are not premises of the written proof.
\end{document}
